\documentclass{report}
\usepackage{amsmath,amssymb}
\setlength{\parindent}{0mm}
\setlength{\parskip}{1em}
\begin{document}
\begin{center}
\rule{6in}{1pt} \
{\large Henricus Bouwmeester \\
{\bf Nonsymmetric multigrid preconditioning for conjugate gradient methods}}

1250 14th St., Suite 600 \\ Denver \\ CO 80202
\\
{\tt henricus.bouwmeester@ucdenver.edu}\\
Andrew Dougherty\\
Andrew V. Knyazev\end{center}

We numerically analyze the possibility of turning off post-smoothing
(relaxation) in geometric multigrid when used as a preconditioner in
conjugate gradient linear and eigenvalue solvers for the 3D Laplacian.
The geometric Semicoarsening Multigrid (SMG) method is provided by the
hypre parallel software package. We solve linear systems using two
variants (standard and flexible) of the preconditioned conjugate gradient
(PCG) and preconditioned steepest descent (PSD) methods. The eigenvalue
problems are solved using the locally optimal block preconditioned
conjugate gradient (LOBPCG) method available in hypre through BLOPEX
software. We observe that turning off the post-smoothing in SMG
dramatically slows down the standard PCG-SMG. For the flexible PCG and
LOBPCG, our numerical results show that post-smoothing can be avoided,
resulting in overall acceleration, due to the high costs of smoothing and
relatively insignificant decrease in convergence speed. We numerically
demonstrate for linear systems that PSD-SMG converges nearly identical to
flexible PCG-SMG if SMG post-smoothing is off. A theoretical
justification is provided.


\end{document}
